\documentclass[../../thesis.tex]{subfiles}

\graphicspath{{Figures/}}

\begin{document}
\chapter*{Introduction}

% Part 1
Modern cosmology is based on the $\Lambda$-CDM model, described in Chapter 1, which predicts the existence of an early phase of exponential expansion of the Universe known as Inflation. This phase plays a central role in our understanding of cosmology, as it is deeply connected to the formation and evolution of the Universe. However, no direct observational evidence supporting this hypothesis has yet been obtained. Consequently, a major goal of current cosmological experiments is to search for such evidence. In particular, efforts are focused on detecting primordial polarisation B-modes patterns in the Cosmic Microwave Background, discussed in Chapter 2. These patterns are expected to be generated by primordial gravitational waves predicted by inflationary theory. At present, only upper limits on the amplitude of this signal have been established from observations.

% Part 2
This absence of detection comes from the multiple difficulties associated with these observations. First, the amplitude of primordial B-modes is theoreticaly expected to be tiny, smaller than former experiment sensitivity \rmk{experiments start to reach the amplitude of these B-modes}. Second, this signal is contaminated by various phenomena: astrophysical foregrounds, gravitational lensing, or atmospheric emissions. The next generation of experiments needs to address these issues by insuring a tight control of systematic effects while allowing separation between primordial B-modes and their contaminations. These are the reasons behind the QUBIC experiment, presented in Chapter 3, called Bolometric Interferometry, an innovative design which combines the high sensitivity of bolometers with unique features of interferometry. In particular, we can also increase our spectral resolution by splitting our physical bandwidth using spectral imaging at the data analysis stage, which allows an improved component separation capability. A focus will be given to spectral imaging in Chapter 4, as it occupies a central place within this thesis. Additionally, interferometry will be a great addition for systematic control, as it allows for self-calibration, a method widely used in radio interferometry. Preliminary work toward this technique has been performed on QUBIC's calibration source and, in particular on its associated GPS, which can be found in Chapter 5. Chapter 6 will be dedicated to the software part of this experiment, as its development has been a large fraction of my time during this PhD.

% Part 3
The main objective of this thesis is the development of map-making methods based on spectral imaging. The first algorithm, called Frequency Map-Making and presented in Chapter 7, reconstructs frequency maps from observational data. Thanks to spectral imaging, multiple maps can be reconstructed within each physical band. A second and improved method, called Component Map-Making, is introduced in Chapter 8. This algorithm aims at directly reconstructing astrophysical component maps, performing map reconstruction and component separation simultaneously, which naturally exploits the spectral imaging capability. Chapter 9 further improves the realism of the simulation pipeline, which is the central piece of the software framework. This includes the modelling of various instrumental and physical parameters, convolution effects, the incorporation of external data, and hyperparameter optimisation. In Chapter 10, an alternative to the forward-modelling formulation of the two algorithms is introduced, providing an approximate solution to the inverse map-making problem using physically-guided neural network approach, guaranteeing the interpretability and allowing noise estimation.

% Part 4
In addition to map reconstruction, separating the primordial CMB signal from astrophysical foregrounds is essential for achieving B-mode detection. Chapter 11 is dedicated to component separation methods and cosmological parameter estimation. These approaches are formulated in spherical harmonic space and rely on the angular power spectra of the CMB and astrophysical emissions. While the previous developments focus on improving map-making and component separation through spectral imaging, they are also applied to a major challenge for ground-based experiments: atmospheric contamination. Traditional strategies rely on data filtering techniques and hardware solutions to reduce atmospheric effects. In contrast, we propose an innovative approach that models the atmosphere as an additional component, evolving in a different domain from astrophysical signals. Chapter 12 details this work, from the development of atmospheric simulations to the mitigation of atmospheric contributions in the data using the Component Map-Making algorithm, and compares the results with those obtained using standard filtering methods. The final chapter, Chapter 13, presents forecasts for the QUBIC instrument. These include comparisons between different sky models and instrumental configurations, as well as between the various algorithms developed throughout this thesis.

\end{document}